\lyseis{Γ1.2}{Συλλογή συνεχών ποσοτικών δεδομένων}
\sol{1} α) $0-10$. β) $10-20$ (το δεξί άκρο της $0-10$ δεν ανήκει σε αυτήν). γ) $10-20$. δ) $20-30$. ε) $0-10$. στ) $30-40$. ζ) $30-40$. η) $20-30$.
\sol{2} α) Πλάτος $10$, κέντρο $35$. β) Πλάτος $10$, κέντρο $155$. γ) Πλάτος $0{,}5$, κέντρο $1{,}25$. δ) Πλάτος $25$, κέντρο $12{,}5$.
\sol{3} Σύνολο: $3+7+6+4=20$. Σχετικές συχνότητες: $\frac{3}{20}=15\%$, $\frac{7}{20}=35\%$, $\frac{6}{20}=30\%$, $\frac{4}{20}=20\%$. Άθροισμα $100\%$.
\sol{4} Οι κλάσεις είναι $5-10$, $10-15$, $15-20$, $20-25$, $25-30$: πέντε κλάσεις. Το $7$ πέφτει στην πρώτη και το $28$ στην τελευταία.
\sol{5} $\frac{4}{40}=10\%$, $\frac{10}{40}=25\%$, $\frac{16}{40}=40\%$, $\frac{8}{40}=20\%$, $\frac{2}{40}=5\%$. Άθροισμα $100\%$.
\sol{6} $8\%$ του $50$ είναι $4$, το $20\%$ είναι $10$, το $36\%$ είναι $18$, το $24\%$ είναι $12$ και το $12\%$ είναι $6$. Άθροισμα: $4+10+18+12+6=50$.
\sol{7} α) Σωστό. β) Σωστό ($\frac{10+20}{2}=15$). γ) Σωστό. δ) Λάθος: ξέρουμε μόνο σε ποια κλάση ανήκει. ε) Σωστό ($\frac{6}{24}=0{,}25$).
\sol{8} Κλάσεις $0-20$: $1$ (το $15$). $20-40$: $4$ ($25$, $30$, $35$, $35$). $40-60$: $6$ ($40$, $45$, $45$, $50$, $55$, $55$). $60-80$: $5$ ($60$, $60$, $65$, $70$, $75$). $80-100$: $4$ ($80$, $85$, $90$, $95$). Σύνολο $20$. Σχετικές συχνότητες: $5\%$, $20\%$, $30\%$, $25\%$, $20\%$.
\sol{9} Κλάσεις $4-8$: $2$ ($6$, $7$), $8-12$: $8$, $12-16$: $10$, $16-20$: $5$. Σύνολο $25$. Σχετικές συχνότητες: $8\%$, $32\%$, $40\%$, $20\%$. Τουλάχιστον $12$ πήραν $10+5=15$ μαθητές, δηλαδή $\frac{15}{25}=60\%$.
\sol{10} Κλάσεις $0-5$: $2$, $5-10$: $6$, $10-15$: $6$, $15-20$: $4$, $20-25$: $2$. Σύνολο $20$ και σχετικές συχνότητες $10\%$, $30\%$, $30\%$, $20\%$, $10\%$. Λιγότερους από $10$ πόντους έβαλε σε $2+6=8$ αγώνες, δηλαδή $\frac{8}{20}=40\%$.
\sol{11} $\frac{12}{80}=15\%$. $25\%$ του $80$ είναι $20$. $\frac{32}{80}=40\%$. $20\%$ του $80$ είναι $16$. Έλεγχος: $12+20+32+16=80$ και $15+25+40+20=100\%$.
\sol{12} α) Κάτω από $6$ kg: $6+12=18$ παιδιά, δηλαδή $\frac{18}{40}=45\%$. β) Τουλάχιστον $8$ kg: $6+2=8$ παιδιά, δηλαδή $\frac{8}{40}=20\%$. γ) Από $4$ ως $8$: $12+14=26$ παιδιά.
\sol{13} Με $12$ κλάσεις κάθε κλάση έχει πολύ λίγες τιμές και ο πίνακας δεν απλοποιεί τίποτα. Με $2$ κλάσεις ο πίνακας χάνει πολλή πληροφορία και δεν φαίνεται πού μαζεύονται οι τιμές. Προτιμότερο είναι να έχουμε περίπου $5$ έως $7$ κλάσεις.
\sol{14} Η συχνότητα της τελευταίας κλάσης είναι $1$ και $\frac{1}{20}=5\%$, όχι $10\%$. Έτσι το άθροισμα βγαίνει $105\%$. Η σωστή τελευταία σχετική συχνότητα είναι $5\%$ και το άθροισμα $100\%$.
\sol{15} Ελεύθερη κατασκευή. Παράδειγμα: $1$, $2$ (κλάση $0-3$), $3$, $4$, $4$, $5$, $5$ (κλάση $3-6$), $6$, $7$, $8$ (κλάση $6-9$). Δεν υπάρχει μόνο μία λύση: υπάρχουν πάρα πολλά σύνολα τιμών με τις ίδιες συχνότητες.
\sol{16} Π.χ. $160$, $161$, $162$, $164$ ή $160$, $163$, $164$, $164$. Ο πίνακας δεν μας λέει ποιες είναι ακριβώς οι τιμές μέσα στην κλάση, μόνο πόσες είναι.
\sol{17} Ελεύθερη κατασκευή. Ο μαθητής πρέπει να έχει $10$ τιμές, κλάσεις πλάτους $2$ χωρίς κενά, και σύνολο συχνοτήτων $10$.
\sol{18} Ελεύθερη κατασκευή. Ο μαθητής πρέπει να έχει $7$ τιμές, $3$ κλάσεις ίσου πλάτους και σύνολο συχνοτήτων $7$.
